\documentclass[conference]{IEEEtran}
\usepackage[utf8]{inputenc}
\usepackage[T1]{fontenc}
\usepackage{cite}
\usepackage{amsmath,amssymb}
\usepackage{graphicx}
\usepackage{booktabs}
\usepackage{url}

\title{Can Deterministic Network Verification Reduce Undetected Faults in LLM‑Generated Configurations?}

\author{
\IEEEauthorblockN{Autonomous AI Researcher}
\IEEEauthorblockA{CNSM 2026 Agentic AI Researcher Track}
}

\begin{document}

\maketitle

\begin{abstract}
We preregistered and executed a paired confirmatory study (study id c1-gnr-vv-2026-0001) to measure whether inserting a deterministic network verifier between an LLM intent\ensuremath{\rightarrow}configuration generator and an emulated execution reduces the proportion of configurations that would execute with operational faults. Using shared\_initial\_candidate semantics (one identical LLM candidate per intent evaluated both without and with verification), we evaluated 72 preregistered paired intents with gpt-5-mini and a canonical deterministic verifier (deterministic\_netops\_validator\_v5). Execution completed with 144 episodes (72 pairs) and 72 model calls. All baseline executions produced no emulator-observed faults (n\_11 = 72, n\_10 = n\_01 = n\_00 = 0). The paired difference in undetected-fault proportion is 0.0 (paired bootstrap 95\% CI [0.0, 0.0], exact McNemar two-sided p = 1.0; bootstrap seed = 1729, resamples = 10000). Under the preregistered shared\_initial\_candidate contract this degenerate confirmatory outcome is expected when identical generated candidates produce no baseline execution faults. We report the preregistration rationale, deterministic execution accounting, representative validator diagnostics, exact inferential outputs, and artifact-grounded recommendations for follow-on confirmatory designs able to detect verifier utility under realistic baseline-error regimes.
\end{abstract}

\section{Introduction}
Generative large language models (LLMs) are increasingly proposed to translate high-level operator intents into device-level network configurations. This intent\ensuremath{\rightarrow}configuration automation can speed change pipelines but risks silently violating reachability, access-control, ordering, or workflow invariants and thereby causing outages if generated text is executed without additional checks. A standard operational mitigation is a generate--verify--repair (GVR) architecture that combines stochastic LLM generation with deterministic verification and, if needed, repair or human review. However, despite many architectural proposals and prototypes, systematic preregistered experiments that measure a verifier's detection performance against executed outcomes under tightly controlled paired designs are scarce.

This paper answers a narrowly scoped, preregistered research question: does inserting a deterministic network verifier between an LLM generator and emulator execution materially reduce the proportion of configurations that execute with operational faults when the identical generated candidate is evaluated both without and with verification? That estimand isolates verifier detection from improvements that could arise from re-sampling, prompt-engineering, or repair-enabled iterations. We executed study c1-gnr-vv-2026-0001 using gpt-5-mini and deterministic\_netops\_validator\_v5 across 72 preregistered paired intents under shared\_initial\_candidate semantics. The run completed with no technical failures, produced a degenerate confirmatory outcome (no baseline emulator faults), and yields a paired difference of 0.0 (95\% CI [0.0,0.0], exact McNemar p = 1.0). Rather than treating this as a null failure, we analyze why it occurred given the preregistered contract and archived artifacts, and we provide concrete, artifact-grounded guidance for designs that can detect verifier utility when baseline error prevalence is non-negligible.

\section{Related Work}
We position this study against three closely related strands: (1) LLM‑based intent\ensuremath{\rightarrow}configuration automation and emulated evaluations; (2) multi-agent and tool‑augmented NetOps frameworks that propose verifier integration; and (3) generate--verify--repair and constrained‑decoding methods from software/code domains that inform possible NetOps adaptations.

LLM intent\ensuremath{\rightarrow}configuration translation. Recent empirical work demonstrates that instruction‑tuned LLMs can map operator intents to device‑level configuration text within constrained vendor dialects and curated testbeds. For example, Nguyen et al. present an intent\ensuremath{\rightarrow}configuration pipeline and evaluate LLM performance on structured configuration tasks, documenting promising task‑level accuracy but highlighting the gap between textual correctness and safe execution in operational environments [1]. Those studies use emulator or synthesized workloads to assess functional correctness of generated artifacts but emphasize that textual matching metrics do not substitute for execution‑level invariants such as reachability or ordering constraints. Our study adopts the same operational framing (emulator execution as a conservative proxy for operational fault) but differs methodologically by preregistering a paired detection estimand under shared\_initial\_candidate semantics to isolate verifier detection capability on identical generated candidates.

Multi-agent and verifier‑integrating NetOps systems. Several prototypes and architectural reports advocate tool‑augmented or multi‑agent LLM systems that incorporate verification steps to reduce regressions and to enable safe automation at scale. Notably, recent multi‑agent intent‑driven frameworks demonstrate how separate agent roles (planner, config generator, validator) can be orchestrated to produce and check configurations before execution; these systems show how verifier signals can be folded into human‑in‑the‑loop or automated repair workflows but typically report system behavior at the demonstration or prototype level rather than in preregistered paired detection trials [2]. Our experiment operationalizes a single, canonical verifier in a strict generate--verify evaluation contract to measure whether the verifier can, by itself and without changing the generated candidate, prevent executed faults that would otherwise occur.

Generate--Verify--Repair and constrained decoding provenance. The generate--verify--repair (GVR) pattern has a substantial methodological pedigree in program repair and structured‑output generation. Architectural proposals and empirical studies in code generation and self‑repair show that deterministic checkers (unit tests, static analyzers) combined with stochastic generation can enable iterative repair loops that converge on functionally correct outputs [3]. Complementary work on constrained decoding and integer‑programming‑based decoding demonstrates that enforcing structural constraints during generation reduces syntactic or semantic violations in structured domains [4]. These results motivate two dimensions of adaptation for NetOps: (a) constraining or biasing generation toward verifier‑satisfiable forms (constrained decoding), and (b) allowing verifier‑triggered repair iterations that change the executed candidate distribution. Importantly, those adaptations change the causal estimand: they do not measure verifier detection on an identical candidate but rather the end‑to‑end benefit of verifier‑enabled generation and repair. The preregistered shared\_initial\_candidate semantics used here intentionally hold the generated candidate fixed across conditions so that any paired difference can be attributed only to verifier detection (not to sampling, constrained decoding, or repair).

Deterministic network verification and policy checking. Deterministic analysis tools developed in the networking literature (header‑space analysis and related formal checkers) provide precise, semantics‑aware invariants for reachability, ACL, and forwarding behavior; these tools are natural verifier candidates for LLM‑generated configs because they can conservatively flag violations that imply operational harm [5]. Prior work on real‑time policy checking and emulated validation shows how formal analyses can be integrated into testbeds to prevent regressions before deployment. Our canonical verifier (deterministic\_netops\_validator\_v5) was used in the same role: apply deterministic, rule‑based checks that compare the generated sequence against manifested workflow constraints (sequence, allowed fields, and protected interfaces) encoded in each task manifest. Where prior literature argues for verifier integration at the architecture level, our contribution is a preregistered paired measurement that quantifies verifier detection power when the generated candidate is held identical across conditions.

Synthesis and gap. Collectively the literature indicates (a) LLMs can produce plausible device commands under constrained settings, (b) multi‑agent and tool‑augmented systems often propose a verifier role, and (c) GVR and constrained decoding strategies can materially change generation outcomes by construction. What is missing---and what this study targets---is a preregistered, paired, artifact‑grounded measurement of verifier detection performance on identical generated candidates: that specific estimand isolates the marginal detection capacity of a verifier decoupled from generation modifications. The degenerate result we observe (zero baseline emulator faults) is thus directly informative: it reveals a design regime---highly constrained synthetic intents and generator behavior---where detection cannot be measured because baseline faults are absent. We use archived traces and validator diagnostics to show how task specification and verifier rule alignment created that regime, and we suggest preregistered alternatives (independent sample arms, repair flows, adversarial task banks) that the literature suggests would expose measurable verifier utility when baseline error prevalence is non-negligible.

\section{Methodology}
Preregistration and estimand. We preregistered study c1-gnr-vv-2026-0001 and froze the design and analysis prior to observing outcomes. The primary estimand is the paired reduction in undetected operational‑fault proportion (paired\_success\_rate\_difference\_guarded\_minus\_baseline). The confirmatory hypothesis H1 targeted an absolute paired reduction of 0.20 with two-sided alpha = 0.05 and power >= 0.80; a sample-size heuristic returned \ensuremath{\approx}63 pairs and we preregistered N = 72 pairs (24 per topology-difficulty stratum) to allow margin for deterministic exclusions.

Shared\_initial\_candidate semantics and experimental contract. The execution\_contract enforced shared\_initial\_candidate semantics: for each intent exactly one initial LLM generation was produced and that identical candidate was evaluated in two conditions. Baseline: execute the shared candidate in an isolated emulator and record emulator fault/no-fault. Guarded: run the canonical deterministic verifier (deterministic\_netops\_validator\_v5) on the same candidate; if the verifier flags violations the guarded outcome is recorded as prevented (no execution); if it does not flag, execute the same candidate in the emulator and record the emulation outcome. This paired structure ensures any observed paired difference arises solely from verifier detection preventing an otherwise-executed fault.

Model, adapter, and verifier configuration. The declared model in execution/model\_configuration.json is gpt-5-mini (provider = openai\_responses, max\_output\_tokens = 2000, maximum\_attempts\_per\_call = 1). The execution record explicitly sets temperature = null; we treat this as provider-default runtime sampling semantics and therefore as a residual, archived sampling-parameter uncertainty. The adapter family used was hosted\_netops\_gvr\_v1. The canonical verifier was deterministic\_netops\_validator\_v5 (validator\_version = "5.0") configured to run the full\_intent\_constraint\_validation rule set. The verifier emits structured validator\_traces per episode that include parsed\_command\_count, required\_command\_count, normalized\_configuration, validation\_before/after final\_state snapshots, violation\_codes, and a difficulty.level tag used for stratification and diagnostics. Repair calls were permitted by the preregistration (maximum\_repair\_calls\_per\_task = 1) but none were applied during the confirmatory episodes (repair\_applied = false in all archived traces).

Task-bank construction and contamination controls. Tasks were deterministically generated by deterministic\_netops\_task\_generator\_v5 (generator\_version = "5.0"). Each task manifest encodes: a natural‑language intent, an ordered machine-structured required\_sequence, allowed\_touched\_fields, initial and expected final states, and a difficulty label (medium/hard/very\_hard). The preregistered contamination procedure performed deterministic exact-match checks against public corpora; analysis/contamination\_summary.csv reports zero exact-match exclusions for confirmatory pairs. Confirmatory stratification (24 pairs per stratum) was preserved as preregistered.

Execution protocol, retries, and deterministic accounting. The missingness\_plan allowed up to two additional retries (3 attempts total) for transient hosted-API errors on generation calls and required full logging of retries and provider events. The run started at 2026-08-31T06:56:23.065707Z and completed at 2026-08-31T07:06:28.665518Z (\ensuremath{\approx}10.1 minutes wall-clock). The execution manifest records planned\_episode\_count = 144, completed\_episode\_count = 144, failed\_episode\_count = 0, and model\_calls\_used = 72 (one shared generation per pair). Provider-call audit (deterministic\_reconciliation.json) shows hosted\_model\_call\_event\_count = 72, completed\_hosted\_model\_call\_event\_count = 72, cache\_miss\_count = 72, and stage\_counts = \{shared\_initial\_generation: 72\}; no retries were consumed for confirmatory pairs.

Scoring, normalization, and binary outcome construction. For each episode the verifier produced a score and validator\_trace. The primary analysis constructs binary indicators per pair as preregistered: baseline\_undetected = 1 iff emulator execution of the shared candidate produced an operational fault; guarded\_undetected = 1 iff the verifier did not flag a violation AND the emulator execution produced an operational fault. Verifier verdicts use normalized\_configuration (deterministic normalization of whitespace and command ordering) and compare parsed\_command\_count against required\_command\_count derived from the task manifest; violation\_codes enumerate rule failures (e.g., ordering\_violation, protected\_interface\_touched). Validator traces for representative episodes show parsed\_command\_count == required\_command\_count and violation\_count = 0 when the candidate satisfied the manifest constraints.

Inferential and bootstrap procedures. The preregistered primary inferential test for the paired binary estimand is an exact McNemar two-sided test when discordant pairs exist. We executed paired\_binary\_analysis\_v1 to produce the 2\ensuremath{\times}2 contingency table (guarded \ensuremath{\times} baseline). Paired bootstrap-percentile 95\% confidence intervals for the paired difference were computed with bootstrap\_seed = 1729 and bootstrap\_resamples = 10,000 (analysis/analysis\_log.jsonl records these parameters). Pre-specified subgroup confirmatory comparisons by topology complexity and policy type were defined and registered with Holm correction, but they are non-informative in the absence of discordant pairs.

Sensitivity, deviations, and diagnostics. Pre-specified sensitivity analyses included: (1) complete\_pair\_primary\_with\_failure\_as\_zero\_sensitivity which treats excluded pairs as failures, and (2) bootstrap diagnostics for verifier true/false detection rates. Because there were no exclusions (analysis/contamination\_summary.csv) and no discordant outcomes (n\_10 = n\_01 = 0), these sensitivity analyses reproduce the degenerate numerical conclusion (paired difference = 0.0). Deterministic reconciliation checks passed: pair accounting and episode accounting are consistent and all deterministic consistency checks passed (deterministic\_reconciliation.json all\_deterministic\_consistency\_checks\_passed = true). All scored episodes and validator\_traces are archived to enable independent reexecution of the same analysis executor.

Artifact grounding. The manifest-level metrics cited above and the verifier configuration (validator\_version = "5.0") are recorded in the archived evidence bundle. Representative per-episode traces under execution/scoring/ include normalized\_configuration, validation\_before/after final\_state, parsed\_command\_count, required\_command\_count, difficulty.level, and repair\_applied flags; these fields are the concrete primitives used by the deterministic scoring rules and the analysis scripts to form binary outcomes. The model configuration file records declared\_model\_version = "gpt-5-mini" and temperature = null; we explicitly report temperature = null as an archived, provider-default runtime sampling semantics parameter that contributes to residual sampling uncertainty in interpretation.

\section{Results}
Execution and manifest summary. The autonomous pipeline completed without technical failures and met the preregistered accounting: planned\_episode\_count = 144; completed\_episode\_count = 144; failed\_episode\_count = 0; model\_calls\_used = 72. The run-level audit (deterministic\_reconciliation.json) records hosted\_model\_call\_event\_count = 72, completed\_hosted\_model\_call\_event\_count = 72, cache\_miss\_count = 72, and stage\_counts = \{shared\_initial\_generation: 72\}. analysis/condition\_summary.csv reports baseline successes = 72, baseline failures = 0 (success\_rate = 1.0) and guarded successes = 72, guarded failures = 0 (success\_rate = 1.0). analysis/paired\_contingency\_table.csv contains the 2\ensuremath{\times}2 counts used for inference.

Paired contingency, inference, and bootstrap diagnostics. Aggregating across the 72 confirmatory pairs produced contingency counts n\_11 = 72, n\_10 = 0, n\_01 = 0, n\_00 = 0 (guarded \ensuremath{\times} baseline). The paired estimate paired\_success\_rate\_difference\_guarded\_minus\_baseline = 0.0. The preregistered exact McNemar two-sided test returns p = 1.0 given zero discordant pairs. The paired bootstrap-percentile 95\% CI (bootstrap\_seed = 1729; bootstrap\_resamples = 10000) is [0.0, 0.0], produced by resampling the observed paired outcomes and recomputing the paired difference; analysis/analysis\_log.jsonl records these bootstrap parameters. These inferential outputs are archived in analysis/paired\_contingency\_table.csv and analysis/condition\_summary.csv.

Per-episode validator traces and representative exemplars. Every episode produced a deterministic validator\_trace (execution/scoring/*.json). Representative extracts from archived per-episode traces show the validator executed full\_intent\_constraint\_validation checks and produced no violation flags across examples: 
- pair-000001 (task-000001): difficulty.level = "medium"; parsed\_command\_count = 4; required\_command\_count = 4; observed\_touched\_field\_count = 3; validator\_version = "5.0"; violation\_count = 0.
- pair-000002 (task-000002): difficulty.level = "hard"; parsed\_command\_count = 2; required\_command\_count = 2; observed\_touched\_field\_count = 2; violation\_count = 0.
- pair-000003 (task-000003): difficulty.level = "hard"; parsed\_command\_count = 6; required\_command\_count = 6; observed\_touched\_field\_count = 4; violation\_count = 0.
The archived manifest lists representative\_tasks with intents, required\_sequence encodings, and responses (execution/responses/task-00000X-*.txt). Across the full confirmatory set the verifier flagged zero violations (analysis/paired\_contingency\_table.csv and analysis/condition\_summary.csv reflect this aggregate result). The per-episode JSONs under execution/scoring/ provide full validator\_before/after state snapshots for reproducibility.

Stratification and subgroup diagnostics. The preregistered stratification (24 pairs per difficulty stratum: medium/hard/very\_hard) was preserved. Because no discordant pairs occurred, the pre-specified subgroup confirmatory comparisons (topology complexity and policy type with Holm correction) produced no informative contrasts; their computation and null outputs are archived but substantively uninformative under a zero-prevalence baseline. The absence of discordance also rendered sensitivity checks (e.g., treating excluded pairs as failures) numerically identical to the main analysis because analysis/contamination\_summary.csv reports zero exact-match exclusions and analysis/missingness\_summary.csv reports complete\_pairs = 72 and zero failed or unscorable episodes.

Sensitivity and exploratory diagnostics executed. We executed the preregistered sensitivity analyses archived in analysis/: (a) complete\_pair\_primary\_with\_failure\_as\_zero\_sensitivity (treating any excluded pairs as failures) --- not applicable because no pairs were excluded; (b) bootstrap CI stability (10k resamples, seed = 1729) --- reproduced CI [0.0,0.0]; (c) exploratory taxonomy of missed faults --- vacuous here because no missed faults were observed. All analysis logs and artifact files for these diagnostics are archived under analysis/ (see analysis/analysis\_log.jsonl, analysis/missingness\_summary.csv, analysis/contamination\_summary.csv).

Artifact-grounded account of the degenerate outcome. The degenerate confirmatory result (no baseline emulator faults) is traceable in archived artifacts to measurable design choices: (1) the task generator produced tightly constrained intents with explicit required\_sequence fields and allowed\_touched\_fields encoded in each task manifest (representative entries in execution/task\_manifest.jsonl), making the desired configuration sequence unambiguous; (2) the model used for initial generation (declared\_model\_version = "gpt-5-mini" in execution/model\_configuration.json) produced candidates that matched the manifest-encoded required sequences across the synthetic corpus; and (3) the verifier's full\_intent\_constraint\_validation rule set directly enforces those same manifest constraints, so verifier expectations and generation targets were highly congruent. Together these factors explain why the shared generated candidates both passed verification and executed without emulator-observed faults.

Reproducibility artifacts and deterministic checks. Key archived artifacts enabling deterministic reconciliation include: execution/raw\_results.jsonl (raw per-episode records), execution/task\_manifest.jsonl (task payloads and required\_sequence encodings), execution/model\_configuration.json (declared\_model\_version = gpt-5-mini; temperature = null; max\_output\_tokens = 2000), analysis/paired\_contingency\_table.csv, analysis/condition\_summary.csv, deterministic\_reconciliation.json, and per-episode validator traces under execution/scoring/. The bootstrap procedure used seed = 1729 and 10,000 resamples; these values are recorded in analysis/analysis\_log.jsonl. Provider-call accounting and model-call counts are recorded in deterministic\_reconciliation.json to allow bit-for-bit reconciliation of run-level totals.

Summary numerical takeaways. Under the preregistered shared\_initial\_candidate semantics and the chosen synthetic task bank the observed baseline undetected-fault proportion = 0/72 = 0.0, guarded undetected-fault proportion = 0/72 = 0.0, paired difference = 0.0, exact McNemar p = 1.0, and paired bootstrap 95\% CI = [0.0,0.0]. All confirmatory diagnostic artifacts and per-episode validator\_traces that substantiate these numbers are archived and referenced above.

\section{Discussion}
Confirmatory interpretation and boundary conditions. The preregistered paired design with shared\_initial\_candidate semantics validly measures verifier detection on identical candidates; because the identical generated candidates produced zero emulator faults the verified paired outcome is a ceiling/null result: the verifier could not reduce executed faults that did not occur. The numeric outputs (n\_11 = 72; paired difference = 0.0; McNemar p = 1.0; bootstrap CI [0.0,0.0]) are direct consequences of the preregistered contract and observed data and therefore correctly reported as a degenerate confirmatory result.

Artifact-grounded causes of the ceiling outcome. Archived artifacts allow us to identify measurable factors that explain the ceiling outcome. Three concrete, artifact-supported points stand out:
1) Task-manifest constraint alignment. The deterministic task generator encoded explicit required\_sequence fields, allowed\_touched\_fields, and workflow\_policies for every confirmatory item (see representative task\_payload entries under manuscript\_evidence\_bundle/representative\_tasks). Validator traces repeatedly show parsed\_command\_count == required\_command\_count for representative episodes (e.g., task-000001: parsed=4 required=4; task-000003: parsed=6 required=6), indicating the generated candidate matched the structured sequence encoded in the task manifest. When generation and manifest constraints are strongly congruent, room for verifier-detectable deviations shrinks.
2) Model--task congruence in a synthetic, constrained bank. The declared model (execution/model\_configuration.json: declared\_model\_version = "gpt-5-mini") produced candidates that matched manifest constraints across the synthetic corpus. Because execution/model\_configuration.json records temperature = null, the run used provider-default sampling semantics; the archived provider-call audit shows hosted\_model\_call\_event\_count = 72 with cache\_miss\_count = 72, indicating fresh calls for each shared generation. Under provider-default sampling and the deterministic, high-clarity prompts used by the generator, the model produced stable, manifest-conformant outputs for the sampled items.
3) Verifier--manifest semantic overlap. The canonical verifier's rule set (deterministic\_netops\_validator\_v5 full\_intent\_constraint\_validation) enforces the same sequence and allowed-field constraints encoded by the task generator. This semantic overlap is visible in per-episode validation\_before/after snapshots archived under execution/scoring/*.json: validator traces show both final\_state snapshots matching expected\_final\_state and violation\_count = 0 across examples. When verifier checks are essentially restatements of the task manifest constraints, and generators respect those constraints, the verifier has limited opportunity to detect faults.

Statistical and design implications. From a statistical point of view the paired estimand requires discordant pairs (n\_10 or n\_01 > 0) to infer detection difference. A zero-prevalence baseline (no executed faults) collapses the paired analysis: the permutation of outcomes offers no information about verifier utility. This outcome does not imply the verifier would always be redundant; it indicates that, for this synthetic bank and sampling semantics, baseline error prevalence was effectively zero. The preregistered power calculation (target paired absolute reduction 0.20) assumed nonzero baseline prevalence; that assumption is violated here and the power argument no longer applies.

What these artifacts imply for follow-on confirmatory designs. The archived evidence suggests concrete, preregisterable modifications that would produce informative, non-degenerate paired comparisons while preserving rigorous controls:
- Independent-generation arm: register independent generator draws per condition (baseline vs guarded) so verifier-mediated repairs or sampling variability can change executed-candidate distributions. This removes the shared\_initial\_candidate constraint and allows measuring downstream benefits of verification+repair or alternative sampling strategies.
- Repair-enabled flows: permit a single verifier-triggered repair call (as allowed by the execution\_contract's transformation\_scope options) and measure whether repaired candidates reduce emulator faults compared to un-repaired executions. The artifact traces already record repair\_applied flags and repair\_prompt\_sha256 placeholders for episodes where repair would be invoked.
- Adversarial or ambiguity-seeded task bank: include items with intentionally underspecified intents, vendor-edge cases, or ordering subtleties where required\_sequence is ambiguous. The deterministic task generator's payload structure supports seeding such ambiguities while preserving provenance checks; analysis/contamination\_summary.csv shows the existing pipeline can detect exact-match contamination and thus preserve preregistration integrity for adversarial items.
- Explicit sampling parameterization: set non-null temperature and record it in execution/model\_configuration.json to quantify sampling-driven variability. The current temperature = null is a residual uncertainty; making sampling explicit will allow controlled exploration of model stochasticity and quantification of verifier value across sampling regimes.
- Multi-model comparisons: include multiple model versions or vendors to measure verifier utility across generators. The deterministic\_reconciliation.json shows the adapter supports logging provider-call audits per model and would enable stratified inference across generators.

Operational tradeoffs and measurable verifier metrics. The artifact vocabulary supports more granular operational metrics than a binary undetected-fault indicator: verifier true-positive rate (proportion of emulator failures flagged), verifier false-positive rate (proportion of verifier flags that would not produce emulator failures), blocking cost (proportion of verifier flags that prevent benign changes), and repair success rate (proportion of repaired candidates that execute without faults). analysis/paired\_contingency\_table.csv and analysis/condition\_summary.csv illustrate how a richer contingency scheme can be recorded. Measuring blocking cost requires permitting verifier-blocked candidates to be inspected or replayed under a repair-enabled flow; the current guarded protocol records prevented outcomes but does not execute prevented candidates, so blocking-cost estimates are not available from this run.

Threats to validity revisited with artifact evidence. Internal validity: the preregistered shared\_initial\_candidate semantics intentionally isolates verifier detection but prevents any verifier-influenced change in executed candidate distributions; this choice explains why a degenerate result may occur in low-error baselines. Construct validity: the emulator-observed fault label is a conservative operational proxy, but emulator-level tests omit traffic-dependent and hardware-specific failure modes; validator traces record final\_state snapshots but cannot capture live-traffic interaction effects. External validity: experiments used a single declared model (gpt-5-mini), a synthetic task bank created by deterministic\_netops\_task\_generator\_v5, and a single deterministic verifier; generalization to other generators, vendor dialects, or production hardware is limited. Conclusion validity: with baseline failure prevalence = 0 the confirmatory paired test lacks the necessary discordance to make inferential claims about verifier usefulness under different baseline regimes.

Practical recommendations grounded in the archive. Based on the archived manifests and traces we recommend that future preregistered studies aiming to quantify verifier utility should:
1) explicitly select or seed a baseline failure prevalence (via adversarial items, ambiguous intents, or historical change examples) so the paired estimand has scope to detect reductions; 2) preregister whether shared\_initial\_candidate semantics or independent-generation semantics will be used, because they answer different operational questions (detection on identical outputs vs. end-to-end improvement under verifier-mediated change); 3) record explicit sampling parameters (non-null temperature) in execution/model\_configuration.json to remove provider-default ambiguity; 4) include repair-enabled arms to measure both detection and net operational improvement; and 5) log verifier-blocking events and, where ethically permissible, execute or simulate blocked candidates to estimate blocking cost. All these recommendations map directly to fields and artifacts present in the current evidence bundle (task manifests, execution/model\_configuration.json, execution/scoring/*.json, and analysis logs) and can be implemented without inventing new tooling.

Concluding perspective. The degenerate confirmatory outcome reported here is an empirically supported datapoint: under the tested constrained, synthetic conditions the generator and manifest constraints produced valid candidates and a canonical verifier did not change executed outcomes. The key scientific value of the run is therefore methodological and procedural: the archived artifacts demonstrate a fully preregistered, deterministic pipeline for measuring verifier detection under tightly controlled semantics, and they make explicit which preregistration choices (shared\_initial\_candidate vs independent generation, sampling parameterization, adversarial seeding, repair permissions) determine whether future confirmatory tests can meaningfully quantify verifier utility.

\section{Limitations}
\begin{itemize}
\item Confirmatory ceiling/null outcome: the baseline generator produced zero emulator faults on the preregistered task set, preventing direct measurement of verifier detection benefit.
\item Shared\_initial\_candidate semantics: the design measures verifier effect on the same generated candidate only and does not assess independent per-condition generation, multi-sample generation, or repairs unless explicitly permitted by the contract.
\item Single-model scope: experiments used one declared model (gpt-5-mini); results do not imply generality across models or versions.
\item Synthetic/emulated tasks: task corpus and emulator executions differ from production devices and live traffic; external validity to production deployments is limited.
\item Sampling-parameter uncertainty: execution/model\_configuration.json shows temperature = null (sampling parameters unset); null indicates provider-default runtime sampling semantics and is a residual uncertainty recorded in the archived configuration.
\item Residual contamination risk: deterministic provenance and exact-match checks were applied, but private training-data contamination of the hosted model cannot be fully ruled out and remains residual uncertainty.
\end{itemize}

\section{Conclusion}
We executed a preregistered paired evaluation (c1-gnr-vv-2026-0001) under shared\_initial\_candidate semantics to test whether a canonical deterministic verifier reduces the proportion of LLM-generated configurations that would execute with operational faults. For the chosen model (gpt-5-mini), verifier (deterministic\_netops\_validator\_v5), and synthetic/emulated task bank, baseline executions produced zero emulator faults and the verifier did not change outcomes (paired difference = 0.0; bootstrap 95\% CI [0.0,0.0]; exact McNemar p = 1.0). This artifact-grounded null/ceiling outcome is internally consistent with the preregistered design: when identical generated candidates produce no executed faults a verifier cannot reduce executed faults under the shared\_initial\_candidate contract. Measuring verifier utility under realistic error regimes requires preregistered contracts that permit independent or multiple generator samples, repair-enabled flows, adversarial task sets, and multi-model comparisons. The archived artifacts (responses, task manifests, validator traces, execution manifest, and analysis logs) provide a reproducible baseline and a template for those follow-on confirmatory designs and power calculations.

\section*{Disclosure Statement}
Autonomy and artifacts: This study was executed autonomously after an autonomy lock. Human role: a Research Director selected the topic family and approved the immutable initial master prompt prior to the autonomy lock; no human manually edited technical manuscript content after the lock. Models and tools: initial generation used gpt-5-mini (declared\_model\_version recorded in execution/model\_configuration.json) orchestrated via the hosted\_netops\_gvr\_v1 adapter; deterministic\_netops\_validator\_v5 (validator\_version = "5.0") served as the canonical verifier; an isolated emulator executed candidate configurations. Preregistration: study c1-gnr-vv-2026-0001 was preregistered and frozen before observing outcomes. Immutable master prompt reference: provenance/master\_prompt.txt (SHA-256: 1872df1e\allowbreak{}1805d2d9\allowbreak{}6940456c\allowbreak{}a016bd66\allowbreak{}5d1d5196\allowbreak{}add77f5a\allowbreak{}cdf1582b\allowbreak{}b39b15ba). Key archived artifacts include execution/raw\_results.jsonl, execution/task\_manifest.jsonl, execution/model\_configuration.json, analysis/paired\_contingency\_table.csv, analysis/condition\_summary.csv, deterministic\_reconciliation.json, and per-episode validator traces under execution/scoring/. The model configuration records temperature = null (provider-default sampling semantics archived in execution/model\_configuration.json). Provider-call audit: hosted\_model\_call\_event\_count = 72. No human manual manuscript editing or content insertion occurred after the autonomy lock.

\begin{thebibliography}{99}
\bibitem{ref1} Nguyen Tu, Sukhyun Nam, James Won‐Ki Hong, "Intent‐Based Network Configuration Using Large Language Models", 2024, 10.1002/nem.2313
\bibitem{ref2} Zhaodong Wang, Samuel Lin, Guanqing Yan, Soudeh Ghorbani, Minlan Yu, Jiawei Zhou, Nathan Hu, Lopa Baruah, Sam Peters, Srikanth Kamath, Jian Yang, Ying Zhang, "Intent-Driven Network Management with Multi-Agent LLMs: The Confucius Framework", 2025, 10.1145/3718958.3750537
\bibitem{ref3} Rafael Cadenas, "Convergent Correctness in Stochastic Code Generation: A Generate-Verify-Repair Architecture for Deterministic Validation of Autonomous AI Coding Agents in Regulated Environments", 2026, 10.2139/ssrn.6754899
\bibitem{ref4} http://citeseerx.ist.psu.edu/viewdoc/summary?doi=10.1.1.300.8659
\end{thebibliography}

\end{document}
